% GENERATED FILE. Edit canonical lessons, then run npm run generate:books.
% canonical-source-hash: fnv1a64:a31dd3b0bd51b125
% canonical-lessons: RU-C57-stone, RU-C57-sand, RU-C57-shore, RU-C57-island, RU-C57-field

\chapter{Stone, Sand, Shore, Island, Field}
\label{ch:ru-stone-sand-shore-island-field}
\hlchaptermodality{57}

\begin{chapteropening}
\textbf{By the end of this chapter:} \emph{I can name a stone, sand, a shore, an island and a field.}

Complete the last lesson of chapter 57: I can name a stone, sand, a shore, an island and a field.
\end{chapteropening}

\section[\texorpdfstring{\ru{камень} (kámen)}{камень (kámen)}]{\ru{камень} (kámen) — a stone}

\label{lesson:RU-C57-stone}



\begin{quote}
Before the new one: say the Russian for a lake, then the Russian for the sea.
\end{quote}

\subsection*{You'll want to know: \ru{камень}}
\textbf{\ru{камень}} (\emph{kámen}) — \textquotedblleft{}a stone\textquotedblright{}.

\begin{cousinweb}
\textbf{\ru{камень}} $\leftarrow$ Proto-Slavic \emph{*kamy}, with the stem \emph{*kamen-} — inherited. Vasmer relates it to Old Norse \emph{hamarr}, \textquotedblleft{}hammer\textquotedblright{}, so English \emph{hammer} is a cousin, beside Lithuanian \emph{akmuõ} and Greek \emph{ákmōn}, \textquotedblleft{}anvil\textquotedblright{}, from a related form.
\end{cousinweb}

\begin{grammarlens}[title={What you've built}]
One more word for this chapter.
\end{grammarlens}

\subsection*{Your turn}
\begin{itemize}
\raggedright
  \item \emph{Say it:} \emph{kámen}
  \item \emph{Say it:} \emph{kámen}, once more
  \item \emph{Say it:} say \emph{kámen}
  \item \emph{Recall:} say the Russian for a lake, then the Russian for the sea, then say \emph{kámen} again
\end{itemize}

\begin{tcolorbox}[breakable,colback=teal!4,colframe=teal!35!black,title={Before you move on}]
What does \ru{камень} mean? (\textquotedblleft{}A stone\textquotedblright{}.) Say it once more.
\end{tcolorbox}
\section[\texorpdfstring{\ru{песок} (pesók)}{песок (pesók)}]{\ru{песок} (pesók) — sand}

\label{lesson:RU-C57-sand}



\begin{quote}
Before the new one: say the Russian for the sea, then the Russian for a stone.
\end{quote}

\subsection*{You'll want to know: \ru{песок}}
\textbf{\ru{песок}} (\emph{pesók}) — \textquotedblleft{}sand\textquotedblright{}.

\begin{cousinweb}
\textbf{\ru{песок}} $\leftarrow$ Proto-Slavic \emph{*pěsŭkŭ} — inherited, a cousin of Sanskrit \emph{pāṃsú-}, \textquotedblleft{}dust, sand\textquotedblright{}, and Avestan \emph{pąsnu-} (Vasmer).
\end{cousinweb}

\begin{grammarlens}[title={What you've built}]
One more word for this chapter.
\end{grammarlens}

\subsection*{Your turn}
\begin{itemize}
\raggedright
  \item \emph{Say it:} \emph{pesók}
  \item \emph{Say it:} \emph{pesók}, once more
  \item \emph{Say it:} say \emph{pesók}
  \item \emph{Recall:} say the Russian for the sea, then the Russian for a stone, then say \emph{pesók} again
\end{itemize}

\begin{tcolorbox}[breakable,colback=teal!4,colframe=teal!35!black,title={Before you move on}]
What does \ru{песок} mean? (\textquotedblleft{}Sand\textquotedblright{}.) Say it once more.
\end{tcolorbox}
\section[\texorpdfstring{\ru{берег} (béreg)}{берег (béreg)}]{\ru{берег} (béreg) — a shore, a bank}

\label{lesson:RU-C57-shore}



\begin{quote}
Before the new one: say the Russian for a stone, then the Russian for sand.
\end{quote}

\subsection*{You'll want to know: \ru{берег}}
\textbf{\ru{берег}} (\emph{béreg}) — \textquotedblleft{}a shore, a bank\textquotedblright{}.

\begin{cousinweb}
\textbf{\ru{берег}} $\leftarrow$ Proto-Slavic \emph{*bergŭ}. Vasmer relates it to Old High German \emph{berg}, \textquotedblleft{}mountain\textquotedblright{}, Avestan \emph{barəzah-}, \textquotedblleft{}height\textquotedblright{}, and Sanskrit \emph{bṛhant-}, \textquotedblleft{}high\textquotedblright{}: a bank is the raised edge. Inherited, that root should have given Slavic \emph{\ru{з}}, not \emph{\ru{г}}, so Berneker and others take the word from Germanic, and Kiparsky from another Indo-European language; Vasmer, judging by its tone and meaning, doubts a loan.
\end{cousinweb}

\begin{grammarlens}[title={What you've built}]
One more word for this chapter.
\end{grammarlens}

\subsection*{Your turn}
\begin{itemize}
\raggedright
  \item \emph{Say it:} \emph{béreg}
  \item \emph{Say it:} \emph{béreg}, once more
  \item \emph{Say it:} say \emph{béreg}
  \item \emph{Recall:} say the Russian for a stone, then the Russian for sand, then say \emph{béreg} again
\end{itemize}

\begin{tcolorbox}[breakable,colback=teal!4,colframe=teal!35!black,title={Before you move on}]
What does \ru{берег} mean? (\textquotedblleft{}A shore, a bank\textquotedblright{}.) Say it once more.
\end{tcolorbox}
\section[\texorpdfstring{\ru{остров} (óstrov)}{остров (óstrov)}]{\ru{остров} (óstrov) — an island}

\label{lesson:RU-C57-island}



\begin{quote}
Before the new one: say the Russian for sand, then the Russian for a shore.
\end{quote}

\subsection*{You'll want to know: \ru{остров}}
\textbf{\ru{остров}} (\emph{óstrov}) — \textquotedblleft{}an island\textquotedblright{}.

\begin{cousinweb}
\textbf{\ru{остров}} $\leftarrow$ Proto-Slavic \emph{*ostrovŭ}, \textquotedblleft{}what the current flows around\textquotedblright{}: \emph{\ru{о}-}, \textquotedblleft{}around\textquotedblright{}, with the root of \emph{\ru{струя}}, \textquotedblleft{}a stream\textquotedblright{} (Vasmer). The root is in Sanskrit \emph{srávati}, \textquotedblleft{}flows\textquotedblright{}, and Greek \emph{rhóos}, \textquotedblleft{}a current\textquotedblright{}.
\end{cousinweb}

\begin{grammarlens}[title={What you've built}]
One more word for this chapter.
\end{grammarlens}

\subsection*{Your turn}
\begin{itemize}
\raggedright
  \item \emph{Say it:} \emph{óstrov}
  \item \emph{Say it:} \emph{óstrov}, once more
  \item \emph{Say it:} say \emph{óstrov}
  \item \emph{Recall:} say the Russian for sand, then the Russian for a shore, then say \emph{óstrov} again
\end{itemize}

\begin{tcolorbox}[breakable,colback=teal!4,colframe=teal!35!black,title={Before you move on}]
What does \ru{остров} mean? (\textquotedblleft{}An island\textquotedblright{}.) Say it once more.
\end{tcolorbox}
\section[\texorpdfstring{\ru{поле} (póle)}{поле (póle)}]{\ru{поле} (póle) — a field}

\label{lesson:RU-C57-field}



\begin{quote}
Before the new one: say the Russian for a shore, then the Russian for an island.
\end{quote}

\subsection*{You'll want to know: \ru{поле}}
\textbf{\ru{поле}} (\emph{póle}) — \textquotedblleft{}a field\textquotedblright{}.

\begin{cousinweb}
\textbf{\ru{поле}} $\leftarrow$ Proto-Slavic \emph{*polje} — inherited, related to Old Russian \emph{polŭ}, \textquotedblleft{}open, empty\textquotedblright{} (Vasmer): a field is open land. Vasmer compares Latin \emph{palam}, \textquotedblleft{}openly\textquotedblright{}, and the \emph{-falen} of \emph{Westphalia}.
\end{cousinweb}

\begin{grammarlens}[title={What you've built}]
One more word for this chapter.
\end{grammarlens}

\subsection*{Your turn}
\begin{itemize}
\raggedright
  \item \emph{Say it:} \emph{póle}
  \item \emph{Say it:} \emph{póle}, once more
  \item \emph{Say it:} say \emph{póle}
  \item \emph{Recall:} say the Russian for a shore, then the Russian for an island, then say \emph{póle} again
\end{itemize}

\begin{tcolorbox}[breakable,colback=teal!4,colframe=teal!35!black,title={Before you move on}]
What does \ru{поле} mean? (\textquotedblleft{}A field\textquotedblright{}.) Say it once more.
\end{tcolorbox}
